%\resizebox{\textwidth}{!}{%
\begin{tikzpicture}

% Layer 1 (bottom): Browser platform
\node[dia/lrbox=13cm] (browser) {%
  \textbf{Браузерна платформа}\\[3pt]
  \textbullet~WebGL --- рендеринг 3D-графіки;\\
  \textbullet~WebRTC --- доступ до камери та мікрофона;\\
  \textbullet~Canvas API --- 2D-графіка та композиція%
};

% Layer 2: Libraries
\node[dia/lrbox=13cm, above=1cm of browser] (libs) {%
  \textbf{JavaScript-бібліотеки}\\[3pt]
  \textbullet~Three.js / A-Frame --- 3D-візуалізація та сцени;\\
  \textbullet~TensorFlow.js / Teachable Machine --- машинне навчання;\\
  \textbullet~WebXR Device API --- доступ до XR-пристроїв%
};

% Layer 3: AR engines
\node[dia/lrbox=13cm, above=1cm of libs] (engines) {%
  \textbf{WebAR-рушії та фреймворки}\\[3pt]
  \textbullet~MindAR --- відстеження зображень, облич;\\
  \textbullet~AR.js --- маркерне та GPS-орієнтоване AR;\\
  \textbullet~8th Wall / Zappar --- комерційні рішення%
};

% Layer 4 (top): Educational application
\node[dia/rbox=13cm, above=1cm of engines] (edu) {%
  \textbf{Освітнє застосування}: імерсивні освітні ресурси (ІОР)%
};

% Arrows between layers
\draw[dia/arrow] (browser.north) -- (libs.south);
\draw[dia/arrow] (libs.north) -- (engines.south);
\draw[dia/arrow] (engines.north) -- (edu.south);

\end{tikzpicture}%
%}
